\section{感想}

ハッシュ表そのものは C 言語の演習で一度実装していたが，同じアルゴリズムを C++ の class として書き直し，さらに実験でハッシュ関数の性質を測ってみると，「動くコード」と「よい設計」の間に多くの論点があることが分かった．特に，加算ハッシュが表の大きさ次第で破綻する様子や，理論式 $1+\alpha/2$ が実測とほぼ完全に一致する様子を自分の実験で確認できたのは面白かった．また，ハッシュ表が 2020 年代でも言語ランタイムの改良対象であり続けていることを文献調査で知り，基礎的なデータ構造の奥行きを感じた．
